\section{Why the changed controller still enumerates refutations}
\label{sec:commands}

This section specializes the already supplied Combined command argument
\cite{combined,cross,new195} to the nine-entry diagonal emitter. The
controller transformation is inherited; its composition with the seven-fold
195-state table is the new construction. Explicit contracts are needed:
changing where modus ponens leaves its result also changes the proof compiler.

Let $p(a,b)=\binom{a+b+1}{2}+b$ and $[a;w]=p(a,w)$. Cantor decoding is
total and satisfies
\[
 \operatorname{hd}[a;w]=a,\qquad
 \operatorname{tl}[a;w]=w,\qquad [0;0]=0.
\]
The last identity must be retained. A numerical stack is not a free list
with a distinct empty constructor. Iterated cons is written
$[a_1,\ldots,a_k;w]$, with $a_1$ on top. Write a command store as
$\langle t,w,d,h\rangle$, where $t$ is the top formula, $w$ the underlying
stack, $d$ an offset, and $h$ the cached formula. Put $I(a,b)=4p(a,b)$.

\subsection{Exact new commands}

The new Shift command is
\[
 S\langle t,w,d,h\rangle=
 \begin{cases}
  \mathsf{halt},&h=1,\\
  \langle h,\operatorname{tl}(w),d,\operatorname{hd}(w)\rangle,&h\ne1.
 \end{cases}
\]
For $y=d+x$, the new modus-ponens command is
\[
 \mathsf{MP}_x\langle t,w,d,h\rangle=
 \begin{cases}
  \langle0,w,0,y\rangle,&I(t,y)=h,\\
  \langle0,w,y,0\rangle,&I(t,y)\ne h.
 \end{cases}                                                    \tag{C1}
\]
Only Shift can halt. Successful MP puts the conclusion in the cache and
clears the offset. It does not push a dummy entry. The old controller
instead returned $\langle y,[0;w],y,0\rangle$ on success and tested
$t+h=1$ at Shift. Thus the two controllers have different store transitions,
even though both recognize refutability of the same calculus.

For $D=d+x$, $H=h+y$, $A=c$, the unchanged AX command has the exact effect
\[
 \mathsf{AX}_{c,x,y}\langle t,w,d,h\rangle
 =\langle K(H,D,c),[q_1,\ldots,q_8,t;w],0,0\rangle,               \tag{C2}
\]
where
\[
 (q_1,\ldots,q_8)
 =(L_6,L_5,L_4,L_3,L_2,L_1,\rho_8(D,A,H,d),U(d,h)).
\]
Here $L_1,\ldots,L_6$ are the six heads in Table~\ref{tab:heads}, and $K$
is the guarded diagonal gate of Section~\ref{sec:pa-calibration}, with zero
on guard failure. All generated values except $U(d,h)$ are implication
codes and have residue zero modulo four; $U(d,h)$ has residue two.
Consequently none of $q_1,\ldots,q_8,K$ is the false code one. These are
numeric residue facts, independent of any consistency assumption.

\subsection{Uniform stack operations}

Commands in a word execute from left to right. For every incoming store,
including $t=1$ and every Cantor zero alias, (C2) gives
\begin{align*}
 \mathsf{AX}_{c,x,y};S^8:
 \langle t,w,d,h\rangle
 &\longmapsto\langle\rho_8(d+x,c,h+y,d),[t;w],0,U(d,h)\rangle,\\
 \mathsf{AX}_{c,x,y};S^9:
 \langle t,w,d,h\rangle
 &\longmapsto\langle U(d,h),w,0,t\rangle.
\end{align*}
The first Shift tests zero and the next eight test generated entries;
none of these nine tests examines the saved top $t$.

Define the following concrete words:
\begin{align*}
 \mathsf{Push}&=\mathsf{AX}_{0,0,0};S^8;\mathsf{MP}_0,\\
 \mathsf{Pop}&=\mathsf{AX}_{0,0,0};S^{10},\\
 \mathsf{Discard}&=\mathsf{MP}_1;\mathsf{AX}_{1,0,0};S^9;\mathsf{MP}_0.
\end{align*}
Their contracts are
\begin{align*}
 \mathsf{Push}\langle t,w,d,h\rangle
 &=\langle0,[t;w],0,0\rangle,\\
 \mathsf{Pop}\langle t,[a;w],d,h\rangle
 &=\begin{cases}
   \mathsf{halt},&t=1,\\
   \langle t,w,0,a\rangle,&t\ne1,
  \end{cases}\\
 \mathsf{Discard}\langle t,w,0,0\rangle&=\langle0,w,0,0\rangle.
\end{align*}
For Push, its last comparison is between an implication and a universal
code, so MP fails with offset zero. Pop's tenth Shift tests the saved top
and, when safe, removes the desired underlying entry. For Discard,
$I(t,1)>0$ forces the first MP to fail, producing offset one. The AX identity
then gives $\langle U(1,0),w,0,0\rangle$, and $I(U(1,0),0)>0$ forces
the final MP to fail with offset zero. Pop is deliberately partial at a
derived contradiction; the proof compiler accounts for this early halt.

Let $\mathsf{Offset}_0$ be empty and let $\mathsf{Offset}_i=\mathsf{MP}_i$
for $i>0$. Starting at $\langle0,w,0,0\rangle$, it produces
$\langle0,w,i,0\rangle$. The word
\[
 \mathsf{Gen}_i=\mathsf{Push};\mathsf{Offset}_i;S;
                    \mathsf{AX}_{i,0,0};S^9
\]
sends $\langle a,w,0,0\rangle$ to $\langle U(i,a),w,0,0\rangle$ for
every $a,w$. AX clears the cache before it could be tested, so the contract
holds even when $a=1$.

\subsection{Axiom selectors and modus ponens}

For $1\le j\le6$, an ordinary head selector is
\[
 \mathsf{Discard};\mathsf{AX}_{A,D,H};S^{j+1};\mathsf{Pop}^{8-j}.
\]
It returns $\langle q_j,w,0,0\rangle$ from
$\langle t,w,0,0\rangle$. The remaining displayed stack entries, ending
in the saved zero, are consumed by the Pops. Each preserves $q_j\ne1$.
For $u\le a$, the restricted propositional head is selected by
\[
 \mathsf{Discard};\mathsf{Offset}_u;\mathsf{AX}_{b,a-u,c};S^8;\mathsf{Pop}.
\]
This returns $\rho_8(a,b,c,u)$. The gate selector
$\mathsf{Discard};\mathsf{AX}_{c,b,a};\mathsf{Pop}^9$
returns $K(a,b,c)$. All preserve the arbitrary underlying stack and leave
both caches zero; their generated tops have residue zero.

The exact MP suffix is
\[
 \mathsf{MPSuffix}_b=S;S;\mathsf{MP}_b;S.
\]
On the store $\langle0,[a,I(a,b),0;w],0,0\rangle$, it halts if $a=1$
or $b=1$, and otherwise returns $\langle b,w,0,0\rangle$. The first
Shift fetches $a$; the second tests it and fetches the major premise.
MP succeeds and puts $b$ in the cache. The final Shift tests the conclusion
and consumes the dummy zero. Without that zero the suffix would consume
an entry of the arbitrary underlying stack. At $a=b=w=0$, total Cantor
decoding proves exactly the same identities.

\begin{lemma}[Uniform proof-word compiler]\label{lem:command-compiler}
For every finite proof tree $p$ in $\mathcal C_\Gamma$, a primitive-recursively
constructed word $W_p$, from every $\langle t,w,0,0\rangle$, either halts
or returns $\langle\operatorname{concl}(p),w,0,0\rangle$.
If $p$ proves one, a word constructed from $p$ halts from every incoming
store, including arbitrary nonzero caches.
\end{lemma}
\begin{proof}
Compile axiom leaves by the preceding selectors and generalization by
$W_{p_0};\mathsf{Gen}_i$. For an MP node with minor premise $p_0$ proving
$a$ and major premise $p_1$ proving $I(a,b)$, use
\[
 W_p=\mathsf{Discard};\mathsf{Push};W_{p_1};\mathsf{Push};
                            W_{p_0};\mathsf{Push};\mathsf{MPSuffix}_b.
\]
The first two words reserve a dummy zero. Induction on $p$ supplies the two
premises above it, with the original stack framed. The suffix has precisely
its required precondition. Any earlier halt remains a halt after extending
the word. This proves the first assertion. If the conclusion is one, use
$\mathsf{AX}_{0,0,0};W_p;\mathsf{Pop}$. The initial AX clears both
caches from every incoming store; if $W_p$ has not already halted, the last
Pop tests its conclusion before inspecting the underlying stack.
\end{proof}

The unchanged scheduler exhaustively enumerates finite command words.
The universal incoming-store property matters: failed earlier trials may
leave a store behind. The successful word need not begin with an empty
stack. The inherited fetch and operand-decoding proofs therefore apply
with the modified command contracts, as their source prefix and all
resolved successors are checked explicitly.

\subsection{Soundness as a finite proof annotation}

Maintain derivations of $t$, $h$, and every canonical stack entry in
$\mathcal C_\Gamma$; impose no derivability condition on $d$. For $w>0$,
$\operatorname{tl}(w)<w$, so iterating the tail reaches zero in at most
$w$ steps. This defines a primitive-recursive finite canonical stack.
The exceptional identity $[0;0]=0$ is handled by the fixed proof
$\rho_8(0,0,0,0)=0$: popping the implicit zero tail always supplies that
proof and requires no infinite annotation.

AX preserves the annotation: it saves $t$, generalizes the old $h$, emits
axioms, and respects $d\le d+x$. MP failure returns zero, while success
applies modus ponens to the annotated $t$ and $h=I(t,d+x)$. Shift moves
annotated entries; a halting Shift exposes an annotated one. These proof
updates are primitive recursive. Induction on a finite command execution
proves that a halt yields a calculus refutation. Conversely,
Lemma~\ref{lem:command-compiler} and exhaustive scheduling turn a calculus
refutation into a finite source halt. Both are PA-verifiable syntactic
constructions. Together with Section~\ref{sec:pa-calibration},
\[
 \PA\vdash \operatorname{Halt}(G_{\rm DC},\mathrm{seed})
 \leftrightarrow
 \exists p\,\operatorname{Prf}_{\PA}(p,\ulcorner0=1\urcorner).     \tag{C3}
\]
Here $G_{\rm DC}$ is the explicitly delivered 3,676-instruction register
program. The next section binds its changed arithmetic loops and its
literal binary-tape implementation.
